\begin{abstract}
% Original, longer abstract retained for reference; shortened version follows.
% The conjugate gradient method (CG) solves a symmetric positive definite (SPD) linear system $A x = b$; it does not, on its own, satisfy the inequality $x \geq 0$. We here develop a \emph{non-negative conjugate gradient} solver for the bound-constrained quadratic
% \[
%   \min_{x}\; \tfrac{1}{2}\,x^\top A x - b^\top x
%   \quad\text{subject to}\quad x \geq 0,\;\; Bx = c,
%   \qquad A \succ 0,
% \]
% that enforces $x \geq 0$ by wrapping CG in an outer active-set loop while retaining CG's $O(\sqrt\kappa)$ Krylov convergence, where $\kappa = \kappa(A)$ is the spectral condition number. The equality system $Bx = c$ is
% present in the equality-augmented variant and absent in the plain bound-constrained and least-squares forms.
%
% We propose the use of a primal-dual active-set method. It solves the unconstrained system over a working set of \emph{free} variables, releases any that come back negative (the primal step), re-admits any bound variable whose reduced gradient is negative (the dual step), and restarts CG on the modified free set each time. These toggles are exactly the principal pivots of the linear complementarity problem $(A,\,{-b})$; because $A$ is a  $P$-matrix, every principal submatrix is invertible and each pivot is well defined. The central result of this paper concerns finite termination by guarding a fast block-pivot path with a least-index Bland fallback. We prove the loop reaches the unique global minimiser in finitely many outer steps with
% \emph{no} non-degeneracy assumption, removing the hypothesis that classical block-principal-pivoting termination arguments require. A synthetic study shows the fallback is never entered on generic data yet genuinely necessary: on an adversarial family of anti-correlated designs the unguarded batch path cycles, and the fallback is what terminates it. The guarantee is robust to inexactness: once CG's residual tolerance is tied to the decision threshold, the inexact loop provably visits the same free sets as the exact one. The guarantee attaches to the outer loop rather than the inner solver, so the active-set method stands as an intermediate result in its own right: the same loop, with the same termination proof, accepts a direct factorisation of each free-set system in place of CG, combining just as readily with
% direct methods when the reduced blocks are small or structured. A loop-free
% projected-gradient reference method, which enforces the constraints simultaneously but forgoes the
% Krylov subspace and so converges at the slower $O(\kappa)$ rate, is analysed alongside it.
%
% Each inner solve is matrix-free: when $A = M^\top M$, the operator $v \mapsto M^\top(M v)$ is
% applied without ever forming $A$, at two products with $M$ per iteration and no $O(n^2)$
% storage, and any regularisation or preconditioner that compresses the spectrum (including a
% ridge/Tikhonov split $A \mapsto (1-\alpha)A + \alpha R^\top R$ that simultaneously secures the
% $P$-matrix property) lowers the iteration count. On synthetic SPD problems the loop reaches
% the planted optimum in at most $\nncgMaxOuter{}$ outer steps, and with a direct free-set solve
% outpaces Lawson--Hanson and an interior-point solver by more than an order of magnitude on
% dense instances. The algorithm is the computational kernel of two companion papers, on
% matrix-free portfolio construction and on random-matrix-theory corrections to the
% mean-variance frontier.

The conjugate gradient method (CG) solves a symmetric positive definite (SPD) linear system $A x = b$ but it does not satisfy the constraint $x \geq 0$. We develop a \emph{non-negative conjugate gradient} solver for the bound-constrained quadratic $\min_x \tfrac{1}{2}x^\top A x - b^\top x$ subject to $x \geq 0$, $Bx = c$, $A \succ 0$, enforcing the constraint by wrapping CG in a primal-dual active-set loop where it solves the unconstrained system on a working set of free variables, releases variables that come back negative, re-admits bound variables with a negative reduced gradient, and restarts CG on the modified free set. These toggles are the principal pivots of the linear complementarity problem $(A, -b)$, well defined because $A$ is a $P$-matrix. Our central result is finite termination if this method is used, meaning that it works in finitely many outer steps and converges to a unique global minimiser with no non-degeneracy assumption, obtained by guarding the fast block-pivot path with a least-index Bland fallback; an adversarial anti-correlated family shows the unguarded path cycles while the fallback terminates it. The guarantee is robust to inexactness once CG's tolerance is tied to the decision threshold, and transfers unchanged to a direct factorisation of each free-set system in place of CG. Each inner solve is matrix-free, applying $A = M^\top M$ as $v \mapsto M^\top(M v)$ without forming $A$; on synthetic SPD problems the loop reaches the planted optimum in at most $\nncgMaxOuter{}$ outer steps and, with a direct free-set solve, outpaces Lawson--Hanson and an interior-point solver by more than an order of magnitude.
\end{abstract}
